\documentclass[11pt]{article}
\usepackage{mystyle}

\title{Rosen --- \S 8.4: Generating Functions\\ \large Selected Exercises}
\author{Alston}
\date{Semester 2, 2026}

\begin{document}
\maketitle

% ======================================================================
\begin{exercise}{17}
    In how many ways can 25 identical donuts be distributed to four
    police officers so that each officer gets at least three but no
    more than seven donuts?
\end{exercise}

% ======================================================================
\begin{exercise}{26}
    \begin{enumerate}[label=(\alph*)]
        \item Show that $1/(1 - x - x^2 - x^3 - x^4 - x^5 - x^6)$ is
        the generating function for the number of ways that the sum
        $n$ can be obtained when a die is rolled repeatedly and the
        order of the rolls matters.
        \item Use part (a) to find the number of ways to roll a total
        of 8 when a die is rolled repeatedly, and the order of the
        rolls matters.
    \end{enumerate}
\end{exercise}

% ======================================================================
\begin{exercise}{39}
    Use generating functions to find an explicit formula for the
    Fibonacci numbers.
\end{exercise}

% ======================================================================
\begin{exercise}{40}
    \begin{enumerate}[label=(\alph*)]
        \item Show that if $n$ is a positive integer, then
        \[
        \binom{2n}{n} = (-4)^n \binom{-1/2}{n}.
        \]
        \item Use the extended binomial theorem and part (a) to show
        that the coefficient of $x^n$ in the expansion of
        $(1-4x)^{-1/2}$ is $\binom{2n}{n}$ for all nonnegative
        integers $n$.
    \end{enumerate}
\end{exercise}

% ======================================================================
\begin{exercise}{41}
    (Calculus required) Let $\{C_n\}$ be the sequence of Catalan
    numbers, that is, the solution to the recurrence relation
    $C_n = \sum_{k=0}^{n-1} C_k C_{n-k-1}$ with $C_0 = C_1 = 1$ (see
    Example 5 in Section 8.1).
    \begin{enumerate}[label=(\alph*)]
        \item Show that if $G(x)$ is the generating function for the
        sequence of Catalan numbers, then $xG(x)^2 - G(x) + 1 = 0$.
        Conclude (using the initial conditions) that
        $G(x) = (1 - \sqrt{1-4x})/(2x)$.
        \item Use Exercise 40 to conclude that
        \[
        G(x) = \sum_{n=0}^{\infty} \frac{1}{n+1} \binom{2n}{n} x^n,
        \]
        so that
        \[
        C_n = \frac{1}{n+1} \binom{2n}{n}.
        \]
        \item Show that $C_n \geq 2^{n-1}$ for all positive integers
        $n$.
    \end{enumerate}
\end{exercise}

% ======================================================================
\begin{exercise}{43}
    Use generating functions to prove Vandermonde's identity:
    \[
    \binom{m+n}{r} = \sum_{k=0}^{r} \binom{m}{r-k} \binom{n}{k},
    \]
    whenever $m$, $n$, and $r$ are nonnegative integers with $r$ not
    exceeding either $m$ or $n$.
    [\emph{Hint}: Look at the coefficient of $x^r$ in both sides of
    $(1+x)^{m+n} = (1+x)^m (1+x)^n$.]
\end{exercise}

% ======================================================================
\begin{exercise}{49}
    A coding system encodes messages using strings of octal (base 8)
    digits. A codeword is considered valid if and only if it contains
    an even number of 7s.
    \begin{enumerate}[label=(\alph*)]
        \item Find a linear nonhomogeneous recurrence relation for
        the number of valid codewords of length $n$. What are the
        initial conditions?
        \item Solve this recurrence relation using Theorem 6 in
        Section 8.2.
        \item Solve this recurrence relation using generating
        functions.
    \end{enumerate}
\end{exercise}

% ======================================================================
\begin{exercise}{56}
    \textit{Recall: $p(n)$, the number of partitions of $n$, is the
    coefficient of $x^n$ in the formal power series expansion of
    $1/((1-x)(1-x^2)(1-x^3)\cdots)$.}

    (Requires calculus) Use the generating function of $p(n)$ to show
    that $p(n) \leq e^{C\sqrt{n}}$ for some constant $C$. [Hardy and
    Ramanujan showed that
    $p(n) \sim e^{\pi\sqrt{2/3}\sqrt{n}}/(4\sqrt{3}n)$, which means
    that the ratio of $p(n)$ and the right-hand side approaches 1 as
    $n$ approaches infinity.]
\end{exercise}

% ======================================================================
\begin{exercise}{57}
    \textit{Recall: the probability generating function for a random
    variable $X$ such that $X(s)$ is a nonnegative integer for all
    $s \in S$ is $G_X(x) = \sum_{k=0}^{\infty} p(X(s)=k)\, x^k$.}

    (Requires calculus) Show that if $G_X$ is the probability
    generating function for such a random variable $X$, then
    \begin{enumerate}[label=(\alph*)]
        \item $G_X(1) = 1$.
        \item $E(X) = G_X'(1)$.
        \item $V(X) = G_X''(1) + G_X'(1) - G_X'(1)^2$.
    \end{enumerate}
\end{exercise}

\end{document}